\chapter{Proposal Feedback and Iteration Plan}
\section{Proposal Feedback Summary}
Project F26-003-Investify was approved with minor modifications. The panel requested clearer contributions beyond leakage, the LLM's user benefit, technical/sentiment explanations, refined problem/model descriptions, explicit metrics, an improved PoC, user transparency, and deeper market understanding.

\section{Proposal Action Register}
The working register preserves the full supplied wording. These identifiers organize project actions; owners and dates require team agreement.
\begin{longtable}{p{0.08\textwidth}p{0.39\textwidth}p{0.36\textwidth}}
\toprule
ID & Required response & Current evidence/status\\
\midrule
PF-01 & Clarify contribution & Revised statement drafted; related-work comparison pending.\\
PF-02 & Explain LLM benefit & Contract/design drafted; real generator pending.\\
PF-03 & Explain technical/sentiment evidence & Complete technical/news context implemented; social/LLM work pending.\\
PF-04 & Refine problem statement & Working revision drafted.\\
PF-05 & Clarify pretrained models & Inventory drafted; social/LLM choices pending.\\
PF-06 & Identify metrics & Matrix and current checks documented.\\
PF-07 & Improve PoC beyond leakage & Corrected verifier/local HTTP demo checked; browser and supervisor demonstration pending.\\
PF-08 & Transparency throughout journey & Data, coverage, and risk limitations exposed; complete failure paths pending.\\
PF-09 & Market understanding & Referenced notes begun; broader research pending.\\
\bottomrule
\end{longtable}

\section{Iteration Strategy}
Each increment supplies a working path, tests, and evidence. The independent code copy retains Git history and uses base commit 1dc24e0. Existing documents and the original checkout remain protected. Upstream access failure is recorded rather than taken to mean no new work exists.

\section{Iteration 1 Plan}
Original commitments cover ingestion, calendar mapping, indicators, news sentiment, social-model selection, the linear baseline, leakage tests/mutation detection, an interface slice, and requirements/design. The first new increment establishes artifact validation and complete evidence presentation. Next work connects and evaluates a real prompted explainer, then addresses failure behavior and risk semantics.

\section{Iteration 2 Plan}
Later work follows the proposal's engine, divergence, deployment, grammar, language-model comparison, runtime, and acceptance milestones. The feedback does not itself require the entire QLoRA programme at this mid.

\section{Work Breakdown and Responsibilities}
Proposal and repo records disagree about harness ownership. The team must reconcile actual contributors and defence responsibilities. Historical contribution is not reassigned or manufactured in this report.

\section{Progress Against Plan}
The selected suite has 136 passes and no skips, including 13 actual Streamlit cases. All 372 records/complete-guided templates, 62 live HTTP checks and 228 separate comparator arithmetic probes validate. A prototype profile/blend, restored viewer, social-study validator/evaluator/view and offline classifier runner are implemented. The source-handover verification is refreshed after each archive change and recorded separately; it excludes external scientific/UI packages. Real LLM/social results, scientific inputs/environment, visual QA and remote CI remain pending.

\section{Delays, Deviations and Recovery Plan}
Repository access, missing data/store, and model runtime constrain full reproduction. The portable artifact path permits tested explanation work while those inputs are restored. The assessment date is unconfirmed. The full mid package is estimated at 45--75 person-hours, subject to data and human-study access.
